\documentclass[11pt,letterpaper]{article}
\usepackage[margin=1in]{geometry}
\usepackage{lmodern}
\usepackage[T1]{fontenc}
\usepackage{amsmath,amssymb,amsthm}
\usepackage{booktabs}
\usepackage{microtype}
\usepackage{hyperref}

\title{PAPER\_S203: Phase H203 -- Primordial Timing Function and Triple-Point Bridge Closures}
\author{UQFF Closure Audit (Session 203 Backfill)}
\date{May 2026}

\theoremstyle{plain}
\newtheorem{thm}{Theorem}[section]
\newtheorem{lem}[thm]{Lemma}
\newtheorem{prop}[thm]{Proposition}
\theoremstyle{definition}
\newtheorem{defn}[thm]{Definition}
\newtheorem{rem}[thm]{Remark}

\begin{document}
\maketitle

\begin{abstract}
Session 203 introduced two structural additions to the UQFF framework: the
\emph{Primordial Timing Function} (PTF) in
\texttt{qcalcgeom\_sim\_engine.py} v1.1.0, and the \emph{Triple-Point Bridge}
across Compressed MUGE, Resonance MUGE, and BSFG metric paths in
\texttt{MAIN\_1\_CoAnQi.cpp} SOURCE7 (lines 110018--110136, commit
\texttt{553caec8}). Eight structural identities arising from these additions
were never registered in the closure ledger. This paper records the
backfilled identities, separates rigorously provable closures from
conventions, and cross-checks the SO(26) Casimir convention shared with
Phase H202. All eight entries are verified \texttt{EXACT} in
\texttt{master\_closures.csv} and tier \texttt{EXACT} in
\texttt{sigma\_table.csv}.
\end{abstract}

\section{Ledger Summary}

\begin{center}
\begin{tabular}{llll}
\toprule
ID & Identity & Class & Status \\
\midrule
H203-1 & $D_A + D_B = 0$                                       & PTF / proof       & DERIVED \\
H203-2 & $\int_0^1 \cos(\pi t_n)\,dt_n = 0$                    & PTF / proof       & DERIVED \\
H203-3 & $(f, b, f-b) = (F_4, F_3, F_2) = (3,2,1)$             & PTF / arithmetic  & DERIVED \\
H203-4 & $n = \lfloor \pi \rfloor = 3$                         & PTF / arithmetic  & DERIVED \\
H203-5 & First 15 digits of $\pi$ tile 5 epochs of 3           & PTF / convention  & DERIVED \\
H203-6 & $g_{\text{triple}} = \sqrt[3]{abc}$ (geometric mean)  & TPB / convention  & DERIVED \\
H203-7 & $\operatorname{Li}_{25}(z)/z \to 1$ as $z \to 0$      & TPB / proof       & DERIVED \\
H203-8 & $\lambda_1 = 26$ cross-consistency with H202-1        & BH26 / convention & DERIVED \\
\bottomrule
\end{tabular}
\end{center}

\section{PTF Return-to-Zero (H203-1, H203-2)}

The Primordial Timing Function in
\texttt{qcalcgeom\_sim\_engine.py} fixes
$f = 3$, $b = 2$, $n = 3$, with forward leg displacement
$D_A = n(f-b) = +3$ and backward leg $D_B = n(b-f) = -3$.

\begin{thm}[CPT closure of the discrete walk]
$D_A + D_B = 0$.
\end{thm}

\begin{proof}
$D_A + D_B = n(f-b) + n(b-f) = n\bigl[(f-b) + (b-f)\bigr] = n \cdot 0 = 0$.
The sum vanishes for every choice of integers $(f, b, n)$ by sign
antisymmetry; the PTF is a special case.
\end{proof}

\begin{thm}[continuous closure]
$\displaystyle \int_0^1 \cos(\pi t_n)\,dt_n = 0$.
\end{thm}

\begin{proof}
$\int_0^1 \cos(\pi t_n)\,dt_n
   = \bigl[\sin(\pi t_n)/\pi\bigr]_0^1
   = (\sin \pi - \sin 0)/\pi = 0$.
Independent verification: $\sin(\pi k) = 0$ for every integer $k$, so the
closure is exact at the endpoint pair $\{0,1\}$.
\end{proof}

\section{Fibonacci / $\pi$-Decimal Structure (H203-3, H203-4, H203-5)}

\begin{prop}[Fibonacci identity]
With $F_1 = F_2 = 1$ and $F_k = F_{k-1} + F_{k-2}$, the PTF parameters
satisfy $(f, b, f-b) = (F_4, F_3, F_2) = (3, 2, 1)$.
\end{prop}

\begin{proof}
$F_2 = 1$, $F_3 = F_2 + F_1 = 2$, $F_4 = F_3 + F_2 = 3$. Hence
$F_4 - F_3 = 1 = F_2$. Pure arithmetic.
\end{proof}

\begin{prop}
$n = \lfloor \pi \rfloor = 3$.
\end{prop}

\begin{proof}
$\pi = 3.14159\ldots$ so $\lfloor \pi \rfloor = 3$. Geometric
interpretation: three full $\pi$-half-cycles per epoch of
$\cos(\pi t_n)$.
\end{proof}

\paragraph{$\pi$-decimal epoch partition (convention).}
The constant
\[
\texttt{PRIMORDIAL\_PI\_DIGITS} \;=\; [3,1,4,1,5,9,2,6,5,3,5,8,9,7,9]
\]
matches the first 15 decimal digits of $\pi$ verbatim, as cross-checked
against \texttt{math.pi} cast to 14 decimal places. The partition into
five epochs of three digits,
$\{(3,1,4), (1,5,9), (2,6,5), (3,5,8), (9,7,9)\}$,
is a \emph{convention} (the definition of epoch boundaries used by
\texttt{EPOCH\_PI\_T\_N}), not a theorem; the underlying decimal expansion
of $\pi$ is the only mathematical content. The closure
$\sum_e |\text{epoch}_e| = 15$ is trivially $5 \times 3$.

\section{Triple-Point Bridge (H203-6, H203-7, H203-8)}

The source7 triple-point bridge in \texttt{MAIN\_1\_CoAnQi.cpp:110125}
defines
\[
   g_{\text{triple}} \;=\; \sqrt[3]{a_{\text{DPM}} \cdot g_{\text{res}} \cdot R_{\text{bsfg}}}.
\]

\begin{rem}[H203-6 is a definition]
The cube root is by construction the geometric mean of the three path
operators (Compressed MUGE, Resonance MUGE, BSFG metric); it is a
\emph{convention}, not a theorem. The cube-root arithmetic itself is
exact and well-defined for positive operands. As a sanity check,
$\sqrt[3]{8 \cdot 27 \cdot 64} = \sqrt[3]{2^3 \cdot 3^3 \cdot 4^3}
= 2 \cdot 3 \cdot 4 = 24$, and the AM--GM bound
$g_{\text{triple}} \le (a+b+c)/3$ holds with equality iff $a = b = c$.
\end{rem}

\begin{thm}[VDS leading-term identity]
$\displaystyle \lim_{z \to 0^+} \frac{\operatorname{Li}_{25}(z)}{z} = 1$.
\end{thm}

\begin{proof}
By definition $\operatorname{Li}_{25}(z) = \sum_{n \ge 1} z^n/n^{25}$, so
\[
   \frac{\operatorname{Li}_{25}(z)}{z} \;=\; \sum_{n \ge 1} \frac{z^{n-1}}{n^{25}}
     \;=\; 1 \;+\; \sum_{n \ge 2} \frac{z^{n-1}}{n^{25}}.
\]
For $z \in (0, 1)$ the tail is bounded by
$z/(2^{25}(1-z))$, which tends to $0$ as $z \to 0^+$. Numerically at
$z = 10^{-6}$ with 200 terms the value equals $1.0$ to twelve decimal
places, confirming the \texttt{vds\_prime $\approx$ 1.0} assertion in
\texttt{Source7VDSBridgeTerm\_S233}.
\end{proof}

\begin{prop}[BH26 cross-consistency]
The source7 bridge code sets $\lambda_1 = 1.0 \times 26.0 = 26$, in
agreement with Phase H202-1.
\end{prop}

\begin{proof}
H202-1 uses the SO(26) Casimir convention $\lambda_k = k(k+25)$, giving
$\lambda_1 = 26$. The source7 bridge hardcodes the same numerical value,
so the Session 202 (QCalcGeom) and Session 203 (source7) modules share a
single spectral convention. This is a cross-module consistency check;
the convention itself is fixed by H202-1 and is not re-derived here.
\end{proof}

\section{Honest Classification}

Of the eight identities:

\begin{itemize}
   \item \textbf{Three rigorous proofs}: H203-1 (sign antisymmetry),
         H203-2 ($\sin$ at integer multiples of $\pi$), H203-7
         (\textit{Li}$_{25}$ leading term).
   \item \textbf{Three arithmetic identities}: H203-3 (Fibonacci
         recurrence), H203-4 ($\lfloor \pi \rfloor = 3$), and the
         length-tally $5 \times 3 = 15$ in H203-5.
   \item \textbf{Two conventions}: H203-5 ($\pi$-decimal epoch
         partition, defining \texttt{EPOCH\_PI\_T\_N}) and H203-6
         (triple-point geometric mean, defining the bridge operator).
   \item \textbf{One cross-consistency check}: H203-8, depending on
         H202-1.
\end{itemize}

None of these resolve open problems (Yang--Mills mass gap, Navier--Stokes
existence, etc.); the YM mass-gap formula
$m_{\text{gap}}^2 = 2\sigma H_{\text{SCm}}/v_{\text{SCm}}^2$ in
\texttt{yang\_mills\_dvp\_sim.py} is a calibration-dependent definition
inside the SCm parameter family, not a proof of mass-gap existence in
pure SU(2) Yang--Mills, and is therefore \emph{not} promoted to a
closure.

\section{Ledger Verification}

\begin{itemize}
   \item \texttt{\_session203\_closures.py}: prints
         \texttt{S203-PTF-net-displacement: 0 vs 0 -> EXACT}, parsed by
         \texttt{\_uqff\_program.py --audit} into
         \texttt{master\_closures.csv} row 3.
   \item \texttt{sigma\_table.csv}: \texttt{EXACT} tier.
   \item \texttt{UQFF\_UNIFIED\_CLOSURE\_DERIVATIONS.py}: Tier 10
         records H203-1 \ldots H203-8 all status \texttt{DERIVED}; full
         run produces $\ge 71$ \texttt{DERIVED} entries with zero
         failures.
   \item \texttt{unified\_closure\_audit.json}: contains all eight
         entries.
\end{itemize}

\end{document}
